\documentclass[10pt,a4paper]{amsart}
\usepackage[margin=19mm]{geometry}
\usepackage{amsmath,amssymb,mathtools}
\usepackage[scr=rsfs,cal=euler]{mathalfa}
\usepackage{xfrac}
\usepackage[unicode,hidelinks,bookmarks=false]{hyperref}
\newcommand{\ie}{i.e.\ }
\newcommand{\cf}{cf.\ }
\newcommand{\resp}{respectively\ }
\newcommand{\Subsection}{\subsection}
\providecommand{\nobreakdash}{\nobreak\hbox{-}}
\setlength{\emergencystretch}{2em}
\multlinegap0pt
\usepackage{bm}
\usepackage{bbm}
\usepackage{dsfont}
\usepackage{xspace}
\usepackage{cancel}
\usepackage{mathtools}
\usepackage[shortlabels]{enumitem}
\usepackage{amsthm}
\makeatletter
\@ifundefined{theorem}{\newtheorem{theorem}{Theorem}}{}
\makeatother
\title[Do quantum linear solvers offer advantage for networks-based]{Do quantum linear solvers offer advantage for networks-based system of linear equations?}
\author{Disha Shetty, Supriyo Dutta, Palak Chawla, Akshaya Jayashankar, Jordi Riu, Jan Nogue, K. Sugisaki, V. S. Prasannaa}
\date{Source: arXiv:2509.00913v4 (CC BY 4.0)}
\begin{document}
\maketitle

\noindent\emph{Source and licence.} Do quantum linear solvers offer advantage for networks-based system of linear equations?. Version arXiv:2509.00913v4 (CC BY 4.0), distributed under the Creative Commons Attribution 4.0 International licence, as recorded in the arXiv licence metadata for this exact version and shown on the abstract page https://arxiv.org/abs/2509.00913v4. Full rights notice: licenses/arxiv-2509.00913-CC-BY-4.0.txt.\par\smallskip

\noindent\textit{Editorial scope.} Counted unit: the Theorem whose source label is \texttt{th:new\_graph} in the article (source0.tex lines 1555--1558, source label \texttt{th:new\_graph}), delivered together with the article's own complete original proof of it (source0.tex lines 1560--1582, one \texttt{proof} environment of 23 lines). No mathematics is written, repaired or re-derived: statement and proof are the article's own text line for line. Required context retained uncounted: none. Every definition, display and label of those lines is retained unchanged. Omitted from this package and not redistributed: 1--1554, 1559--1559, 1583--2180. Nothing inside a retained excerpt is omitted or altered: every retained body is delivered exactly as the article's lines are written, so no omission receipt of any kind accompanies this package. Citations: the retained excerpts carry no citation command at all, so no bibliography entry of the article is transcribed and no reference is redistributed. Reference adaptation: none was needed and none was made; every reference inside the delivered unit resolves inside it, so no unresolved reference or citation is produced. Printed slips kept literal: no printed word, symbol, hypothesis or number of the counted statement or of its proof was altered. Layout and class: the article's own class is \texttt{revtex4-2} with packages including float, graphicx, dcolumn, bm, babel, inputenc, amsthm, mathtools, physics, algorithm, algpseudocode, amsmath, xcolor, geometry; the wrapper is the lane's pinned \texttt{amsart} editorial wrapper at 10pt on A4, which restates the theorem-like environments the retained text needs and adds the article's own macro definitions that the retained text needs (none), with extra wrapper packages (bm, bbm, dsfont, xspace, cancel, mathtools, enumitem). The delivered numbering is the unit's own. Assets: no retained span contains a figure environment, a graphics-inclusion command, a PDF-page inclusion, a table or a TikZ picture; no image file is copied into this package, the package declares no asset dependency and no image file is redistributed with it. Licence: version arXiv:2509.00913v4 of this article is distributed under the Creative Commons Attribution 4.0 International licence, as recorded in the arXiv licence metadata for this exact version and shown on the abstract page https://arxiv.org/abs/2509.00913v4. A full rights notice is delivered at licenses/arxiv-2509.00913-CC-BY-4.0.txt. Characters: every retained excerpt is pure ASCII, so the delivered engine, XeLaTeX with its default Latin Modern text font, reports no missing character.
\begin{theorem}\label{th:new_graph}
Let $G = (V(G), E(G))$ be any graph with the set of vertices $V(G) = \{v_i:i=1,\ 2,\ 3, \cdots, N\}$ and set of edges $E(G)$. Let $H = (V(H), E(H))$ be a graph such that $V(H) = V(G) \cup \{v_{(N + 1)}, v_{(N + 2)}\}$ and set of edges $E(H) = E(G) \cup \{(v_{(N + 1)}, u_1), (v_{(N + 1)}, u_2), \cdots, (v_{(N + 1)}, u_k)\} \cup \{(v_{(N + 2)}, u_1), (v_{(N + 2)}, u_2), \cdots, (v_{(N + 2)}, u_k)\}$,
where $u_1, u_2, \dots u_k \in V(G)$ in some order. Then, the Laplacian matrix $L(H)$ has eigenvalue $k$ with an eigenvector $\vec{b} = (\underbrace{0,0,\ldots,0}_{N \text{ times}}, 1, -1)^{\mathsf{T}}$.
\end{theorem}

\begin{proof}
    We know that $L(H) = D(H) - Q(H)$, where $D(H)$ and $Q(H)$ are the degree matrix and the adjacency matrix of the graph $H$, respectively. Denote $L(H) = (l_{ij})_{(N + 2) \times (N + 2)}$ where
	$$l_{ij} = \begin{cases} d_i^{(H)} &\text{if}~ i = j; \\ -1 &\text{if}~ (i, j) \in E(H); \\ 0 &\text{otherwise};
	\end{cases}$$
	where $d_i^{(H)}$ is the degree of vertex $v_i$ in the graph $H$. Now multiply $L(H)$ with the vector $\vec{b}$. The $i$-th row of $L(H)\times \vec{b}$ is
	\begin{equation}
		(L(H) \vec{b})_{i \bullet} = \sum_{j = 1}^{N + 2} l_{ij} b_j,
	\end{equation}
	where $i = 1, 2, \cdots, (N + 2)$. Let $b_1 = b_2 = \cdots = b_N = 0$. Now, consider the following cases: 
	\begin{enumerate}[label = {Case \arabic*:}]
		\item 
		Let $v_i$ be a vertex other than $v_1, v_2, \cdots, v_k, v_{(N + 1)}$, and $v_{(N + 2)}$. Now, $l_{i (N + 1)} = l_{i (N + 2)} = 0$. Also, $b_1 = b_2 = \cdots = b_N = 0$. Therefore, $(L(H) \vec{b})_{i \bullet} = 0$.
		\item 
		Let $u_q$ be any one of $u_1, u_2, \cdots, u_k$. Note that the rows of $L(H)$ corresponding to $u_q$ always have two non-zero entries $l_{q (N + 1)} = l_{q (N + 2)} = -1$. As $b_1 = b_2 = \cdots = b_N = 0$, we have $\sum_{j = 1}^{N} l_{q j} b_j = 0$. As $b_{(N + 1)}$ and $b_{(N + 2)}$ have opposite signs, we have $l_{q (N + 1)} b_{(N + 1)} + l_{q (N + 2)} b_{(N + 2)} = 0$. 
		\item 
		Let $i = (N + 1)$ or $i = (N + 2)$. Note that the degrees of $(N + 1)$ and $(N + 2)$ are $d_{(N + 1)}^{(H)} = d_{(N + 2)}^{(H)} = k$. Therefore $l_{(N + 1)(N + 1)} = l_{(N + 2)(N + 2)} = k$. Now, $(L(H) b)_{(N + 1), \bullet} = k$ and $(L(H) b)_{(N + 2), \bullet} = -k$. 
	\end{enumerate}
	Combining all the cases, we observe that 
	\begin{equation}
		L(H)\vec{b} = (\underbrace{0,0,\ldots,0}_{n \text{ times}}, k, -k)^{\mathsf{T}} = k \vec{b}. 
	\end{equation}
	Therefore, $L(H)$ has an eigenvalue $k$ with eigenvector $\vec{b}$.
\end{proof}

\end{document}
